\documentclass[aspectratio=169, 10pt]{beamer}
\usetheme{metropolis}

% --- Edinburgh colours ----------------------------------------
\definecolor{UoERed}{HTML}{7A2318}
\definecolor{UoEGold}{HTML}{B8860B}
\definecolor{CodeBG}{HTML}{F7F3ED}
\definecolor{CodeFrame}{HTML}{D5C9B1}

\setbeamercolor{palette primary}{bg=UoERed, fg=white}
\setbeamercolor{frametitle}{bg=UoERed, fg=white}
\setbeamercolor{title separator}{fg=UoEGold}
\setbeamercolor{progress bar}{fg=UoEGold, bg=UoEGold!20}
\setbeamercolor{alerted text}{fg=UoERed}
\setbeamercolor{block title}{bg=UoERed!15, fg=UoERed}
\setbeamercolor{block body}{bg=UoERed!5}

% --- Packages -------------------------------------------------
\usepackage[T1]{fontenc}
\usepackage{lmodern}
\usepackage{amsmath, amssymb, mathtools}
\usepackage{listings}
\usepackage{graphicx, booktabs}
\usepackage{tikz}
\usetikzlibrary{arrows.meta, positioning, calc, decorations.pathreplacing}
\usepackage{hyperref}
\hypersetup{colorlinks=true, linkcolor=UoERed, urlcolor=UoERed!70!black}
\setbeamertemplate{navigation symbols}{}

\lstset{
  language=Python,
  basicstyle=\ttfamily\scriptsize,
  keywordstyle=\color{UoERed}\bfseries,
  commentstyle=\color{gray}\itshape,
  stringstyle=\color{UoEGold},
  backgroundcolor=\color{CodeBG},
  frame=single,
  rulecolor=\color{CodeFrame},
  numbers=none, breaklines=true, showstringspaces=false, tabsize=4,
  xleftmargin=4pt, xrightmargin=4pt, aboveskip=6pt, belowskip=4pt,
}

\AtBeginSection[]{
  \begin{frame}
    \vfill\centering
    \usebeamerfont{title}\insertsectionhead\par
    \vfill
  \end{frame}
}

\title{Week 5 --- Deep Equilibrium Networks II}
\subtitle{Stochastic Models \& Gauss--Hermite Quadrature}
\author{Dr Juan Zurita}
\institute{Deep Learning for Macroeconomics\\
The University of Edinburgh~$\cdot$~School of Economics}
\date{}

\begin{document}
\maketitle

\begin{frame}{Adding Uncertainty}
Real economies face shocks: $y_t = e^{z_t} k_t^\alpha$, where $z_t$ follows an AR(1).

\bigskip
The Euler equation now involves an \alert{expectation}:

$$u'(c_t) = \beta\,\mathbb{E}_t\!\left[u'(c_{t+1})\,(1 + r_{t+1})\right]$$

\bigskip
How do we compute $\mathbb{E}_t[\cdot]$ inside a neural-network training loop?
\end{frame}

\begin{frame}{Gauss--Hermite Quadrature}
Approximate the expectation with a weighted sum:

$$\mathbb{E}[g(\varepsilon)] \approx \sum_{j=1}^{J} \omega_j\,g(\varepsilon_j)$$

\bigskip
\begin{itemize}
\item $\varepsilon_j$: quadrature nodes (roots of Hermite polynomials)
\item $\omega_j$: quadrature weights
\item With $J = 5$--$7$ nodes, accuracy is excellent for smooth integrands
\end{itemize}

\bigskip
This integrates seamlessly with PyTorch --- each node is a forward pass through the network.
\end{frame}

\begin{frame}{Stochastic Brock--Mirman}
State: $(k_t, z_t)$. Policy: $c_t = \mathcal{N}_{\boldsymbol{\theta}}(k_t, z_t)$.

\bigskip
Residual with quadrature:

$$R(k,z; \boldsymbol{\theta}) = 1 - \beta \sum_j \omega_j \frac{\alpha (k')^{\alpha-1} e^{z'_j}}{c(k', z'_j; \boldsymbol{\theta}) / c(k,z; \boldsymbol{\theta})}$$

\bigskip
Now the network takes \textbf{two} inputs --- this is where the grid-free advantage starts to matter.
\end{frame}

\begin{frame}{Loss Function Variants}
\begin{center}
\begin{tabular}{lll}
\toprule
Loss & Formula & Properties \\
\midrule
MSE & $R^2$ & Standard, sensitive to outliers \\
Huber & smooth $L_1$ & Robust to large residuals \\
Log-cosh & $\ln(\cosh R)$ & Smooth Huber approximation \\
\bottomrule
\end{tabular}
\end{center}

\bigskip
In practice, \textbf{MSE} works well for most economic models. Switch to Huber if training is unstable.
\end{frame}

\begin{frame}{Summary}
\begin{itemize}
\item Stochastic models require numerical integration of expectations
\item Gauss--Hermite quadrature is accurate and differentiable
\item The DEQN approach extends naturally to stochastic settings
\item With two state variables, we already outperform naive grids
\end{itemize}

\bigskip
\textbf{Next week:} occasionally binding constraints --- the Fischer--Burmeister trick.
\end{frame}

\end{document}
